\documentclass[../DoAn.tex]{subfiles}
\begin{document}

Chương này trình bày chi tiết về các công nghệ, nền tảng phát triển và công cụ quản lý phiên bản được sử dụng trong dự án trò chơi "Dream Witch". Việc lựa chọn công nghệ được thực hiện dựa trên cơ sở phân tích kỹ lưỡng các yêu cầu kỹ thuật và lối chơi được đặc tả ở Chương 2, nhằm đưa ra giải pháp tối ưu cho hiệu năng và tiến độ phát triển của dự án.

\section{Hệ thống camera thông minh (Cinemachine)}
\label{section:tech:3.1}

Để mang lại trải nghiệm hình ảnh điện ảnh và bám sát các chuyển động tốc độ cao của nhân vật chính (như nhảy kép, lướt), dự án đã tích hợp hệ thống camera thông minh Cinemachine. Khác với hệ thống camera tĩnh hoặc lập trình thủ công truyền thống, Cinemachine cung cấp các công cụ cấu hình theo mô hình đạo diễn ảo (Virtual Camera) cho phép thiết lập và chuyển đổi nhiều góc máy một cách linh hoạt. Do đặc thù là trò chơi 2D, camera chính của hệ thống được thiết lập ở chế độ trực giao (Orthographic Camera), loại bỏ hiện tượng biến dạng phối cảnh (perspective distortion), giúp duy trì tỷ lệ kích thước chuẩn xác của Sprites bất kể khoảng cách.

So với việc tự lập trình bằng mã nguồn C\# cho phép nội suy vị trí camera, Cinemachine xử lý mượt mà bài toán theo dấu mục tiêu thông qua thành phần Framing Transposer. Sự kết hợp giữa vùng chết (Dead Zone) và giới hạn mềm (Soft Zone) giúp camera không liên tục bị giật (jitter) theo từng bước di chuyển nhỏ của nhân vật. Đặc biệt, hệ thống tích hợp thành phần mở rộng Cinemachine Confiner 2D kết nối trực tiếp với Collider 2D của từng phòng chơi (Room Controller), đảm bảo khung hình camera luôn bị giới hạn trong khu vực cho phép và không bao giờ lọt ra ngoài không gian hổng của bản đồ. Cinemachine cũng cung cấp sẵn module Impulse để dễ dàng tạo hiệu ứng rung lắc màn hình (Camera Shake) khi nhận sát thương hoặc khi Boss đập đất.

% NOTE: Sinh viên có thể chèn hình chụp thực tế thiết lập Cinemachine tại đây
\begin{figure}[H]
    \centering
    \includegraphics[width=0.9\linewidth,height=0.25\textheight,keepaspectratio]{Hinhve/CinemachineCamera.png}
    \caption{Thiết lập thông số kỹ thuật cho Virtual Camera trong giao diện Cinemachine}
    \label{fig:cinemachine_editor}
\end{figure}

\section{Hệ thống kịch bản và hội thoại đa nhánh (Yarn Spinner)}
\label{section:tech:3.2}

Để hiện thực hóa yêu cầu kể chuyện và hướng khám phá của người chơi (Mục \ref{section:2.2}), trò chơi cần tích hợp một hệ thống quản lý hội thoại đa nhánh với các NPC (Shopkeeper, Facility Owner). Đề tài lựa chọn Yarn Spinner \cite{yarnspinner} để quản lý hội thoại động và liên kết các sự kiện trong game. Yarn Spinner cho phép người chơi lựa chọn các phương án trả lời khác nhau, hỗ trợ liên kết trực tiếp giữa các câu lệnh trong văn bản hội thoại (Yarn Script) với các hàm C\# của hệ thống game (ví dụ kích hoạt giao diện mua sắm khi nói chuyện với Shopkeeper thông qua \texttt{YarnCommand}), đồng thời đồng bộ hóa các biến số hội thoại với hệ thống lưu trữ dữ liệu tổng của trò chơi.

So với việc tự lập trình bộ phân giải JSON/XML (tốn công sức phát triển, khó bảo trì khi kịch bản dài) hay sử dụng ngôn ngữ Ink (khó gọi ngược các hàm sự kiện UI của Unity do tiếp cận theo hướng dòng chảy văn bản liên tục), Yarn Spinner là lựa chọn tối ưu nhờ cấu trúc viết kịch bản dạng nút (node-based) trực quan. Cú pháp viết script đơn giản cùng tính năng kiểm tra lỗi cú pháp trực tiếp trong Unity Editor giúp giảm thiểu tối đa thời gian thiết kế và phát triển kịch bản hội thoại.

% NOTE: Sinh viên có thể chèn hình chụp thực tế giao diện thiết kế nút hội thoại (Yarn Editor) tại đây
\begin{figure}[H]
    \centering
    \includegraphics[width=0.9\linewidth,height=0.28\textheight,keepaspectratio]{Hinhve/YarnEditor.png}
    \caption{Giao diện thiết kế các nút hội thoại (Node) trong Yarn Editor}
    \label{fig:yarn_editor_screenshot}
\end{figure}

\section{Công nghệ hoạt ảnh 2D dạng xương (Spine Animation)}
\label{section:tech:3.3}

Để hiện thực hóa trải nghiệm game hành động mượt mà với hoạt ảnh Boss linh hoạt phức tạp (Mục \ref{section:2.2}), đề tài lựa chọn công nghệ Spine Animation \cite{esotericspine}. Spine giải quyết bài toán tài nguyên và hiệu năng nhờ kỹ thuật biến dạng lưới khung xương (Skeletal Mesh Deformation) giúp tạo ra các chuyển động trượt, nén mượt mà mà không phụ thuộc vào số lượng hình ảnh vẽ tay lớn. Công nghệ này cũng tự động hòa trộn hoạt ảnh (Animation Blending) và cho phép gắn các sự kiện vào mốc thời gian hoạt ảnh (Keyframe Event) để đồng bộ hóa các hitbox tấn công cận chiến chính xác.

So với phương pháp Sprite Sheets truyền thống vẽ tay từng khung hình (tiêu tốn lượng VRAM khổng lồ khi nhân vật có nhiều biến thể hành động) hay bộ công cụ Unity Sprite Skinning mặc định (thiếu các tính năng điều khiển xương động nâng cao như IK, Path Constraints), Spine cung cấp runtime tối ưu hóa vượt trội cho Unity. Nó giúp giảm tối đa dung lượng bộ nhớ kết cấu bằng việc tái sử dụng một bộ xương cho nhiều giao diện (Skins) và nâng cao trải nghiệm thị giác nhờ hiệu ứng giả lập 3D sống động cho các thực thể lớn.

% NOTE: Sinh viên có thể chèn hình chụp thực tế giao diện thiết kế xương nhân vật trong phần mềm Spine tại đây
\begin{figure}[H]
    \centering
    \includegraphics[width=0.9\linewidth,height=0.28\textheight,keepaspectratio]{Hinhve/SpineEditor.png}
    \caption{Khung xương của nhân vật thiết kế trong phần mềm Spine}
    \label{fig:spine_editor_screenshot}
\end{figure}

\section{Giải thuật tìm đường thông minh (A* Pathfinding)}
\label{section:tech:3.4}

Các kẻ địch bay trong game đòi hỏi khả năng định tuyến di chuyển và săn đuổi người chơi một cách tự nhiên. Hệ thống sử dụng thư viện A* Pathfinding Project \cite{arongranbergastar} phát triển trên giải thuật A* để giải quyết vấn đề tìm đường thông minh của thực thể AI. Hệ thống hỗ trợ tự động quét Tilemap để tạo lưới đồ thị di chuyển (Grid Graph), cập nhật đường đi động theo thời gian thực (khi cửa đóng, bẫy kích hoạt) và làm mịn đường dẫn (Path Smoothing) giúp chuyển động của các quái vật bay không bị gấp khúc cơ học.

Các giải pháp thay thế như Unity NavMesh 2D (chủ yếu thiết kế cho 3D, hoạt động trên 2D thiếu ổn định, không tối ưu cho địa hình platformer phức tạp) hay tự cài đặt BFS/Dijkstra (khó tối ưu CPU overhead khi tính toán cho hàng chục kẻ địch đồng thời) đều bộc lộ hạn chế. Thư viện A* Pathfinding Project được chọn nhờ tích hợp sâu với Unity Job System và Burst Compiler, cho phép tính toán đa luồng độ trễ cực thấp. Ngoài ra, thư viện còn hỗ trợ gán trọng số (penalties) cho các vùng nguy hại (gai, axit) để AI chủ động tránh né.

% NOTE: Sinh viên có thể chèn hình chụp thực tế lưới quét tìm đường A* (A* Grid Graph) hiển thị trong Unity Scene View tại đây
\begin{figure}[H]
    \centering
    \includegraphics[width=0.9\linewidth,height=0.28\textheight,keepaspectratio]{Hinhve/AstarGridGraph.png}
    \caption{Lưới đồ thị di chuyển (Grid Graph) hiển thị trực quan dạng Gizmos trong Unity Scene View}
    \label{fig:astar_grid_graph}
\end{figure}


\end{document}